\documentclass[UTF8,fontset=none,11pt]{ctexart}

\usepackage[a4paper,margin=23mm,headheight=16pt]{geometry}
\usepackage{fontspec}
\setmainfont{texgyrepagella-regular.otf}[BoldFont=texgyrepagella-bold.otf,ItalicFont=texgyrepagella-italic.otf,BoldItalicFont=texgyrepagella-bolditalic.otf]
\setsansfont{texgyreheros-regular.otf}[BoldFont=texgyreheros-bold.otf]
\setmonofont{lmmono10-regular.otf}[Scale=0.88,BoldFont=lmmonolt10-bold.otf]
\usepackage{amsmath,amssymb,booktabs,longtable,array,tabularx}
\usepackage{xcolor,listings,enumitem,fancyhdr,tikz}
\usetikzlibrary{arrows.meta,positioning}
\usepackage{hyperref,xurl}
\definecolor{ink}{HTML}{17324D}
\definecolor{teal}{HTML}{087F8C}
\definecolor{pale}{HTML}{F2F6F9}
\hypersetup{colorlinks=true,linkcolor=ink,urlcolor=teal,bookmarksnumbered=true}
\urlstyle{tt}
\lstset{basicstyle=\ttfamily\small,breaklines=true,columns=fullflexible,
 keepspaces=true,showstringspaces=false,frame=single,rulecolor=\color{teal!25},
 backgroundcolor=\color{pale},xleftmargin=3pt,xrightmargin=3pt,
 aboveskip=9pt,belowskip=9pt,commentstyle=\color{teal},tabsize=4}
\setlist{nosep,leftmargin=1.5em,topsep=5pt}
\setlength{\parskip}{5pt}
\setlength{\parindent}{0pt}
\setlength{\emergencystretch}{3em}
\renewcommand{\arraystretch}{1.25}
\pagestyle{fancy}
\fancyhf{}
\fancyhead[L]{\small\sffamily\color{ink}QUANTLAB}
\fancyhead[R]{\small\sffamily\color{ink}\nouppercase{\leftmark}}
\fancyfoot[C]{\small\thepage}
\renewcommand{\sectionmark}[1]{\markboth{\thesection\quad #1}{}}
\newcommand{\code}[1]{\nolinkurl{#1}}
\newcommand{\note}[1]{\par\smallskip\noindent\colorbox{pale}{\parbox{\dimexpr\linewidth-2\fboxsep\relax}{#1}}\par\smallskip}
\newcommand{\chapterpage}{\par}
\usepackage{needspace}
\let\originalsection\section
\renewcommand{\section}{\Needspace{9\baselineskip}\originalsection}
\let\originalsubsection\subsection
\renewcommand{\subsection}{\Needspace{8\baselineskip}\originalsubsection}
\widowpenalty=10000
\clubpenalty=10000

\IfFontExistsTF{STSong}{\setCJKmainfont{STSong}[BoldFont=Heiti SC]\setCJKsansfont{Heiti SC}\setCJKmonofont{STSong}}{\setCJKmainfont{FandolSong-Regular.otf}[BoldFont=FandolHei-Regular.otf]\setCJKsansfont{FandolHei-Regular.otf}\setCJKmonofont{FandolSong-Regular.otf}}
\hypersetup{pdftitle={QuantLab 项目、使用与维护手册},pdfauthor={QuantLab}}

\begin{document}
\begin{titlepage}
\color{ink}
{\sffamily\large LOCAL RESEARCH / AUDITABLE BACKTESTING\par}
\vspace{28mm}
{\sffamily\fontsize{42}{48}\selectfont QuantLab\par}
\vspace{8mm}
{\Huge 美股量化研究系统\par}
\vspace{5mm}
{\LARGE 项目、使用与维护手册\par}
\vspace{12mm}
{\large 中文版\par}
\vfill
\rule{\linewidth}{1pt}
\vspace{5mm}
从策略想法到可审计的历史实验：数据、因子、策略、成交、账户、指标、报告与桌面界面。

\vspace{8mm}
版本：项目 2026.10.22\quad 核对日期：2026 年 10 月 22 日

适用目录：\path{/Users/yusa/ChatGPTWorkSpace/QuantLab}

本手册根据当前本地代码编写。它区分已实现功能、独立研究脚本和未来扩展；不将历史回测表现作为未来收益承诺。
\end{titlepage}
\tableofcontents
\clearpage

\section{项目是什么，以及适合谁}
QuantLab 是一个在本机运行的美股日线研究与回测系统。用户提供价格数据、策略代码和运行配置，系统按明确的时间顺序模拟历史交易，维护现金与持仓，计算绩效，并保存可以复核的完整实验档案。它的重点是正确性、可审计性和可复现性，而不是自动发现收益最高的策略。

项目包含三个层次：Python 回测引擎是计算核心；命令行与 Tauri 桌面 App 是两种操作入口；\path{researchProject/} 保存专题研究及其独立计算材料。桌面界面不会替换引擎的交易规则，研究专题也不等同于可以在策略菜单中直接运行的模块。

\subsection{能够完成什么}
\begin{itemize}
\item 读取本地 CSV/Parquet 日线，检查价格、股票代码、交易日和基准数据。
\item 执行用户编写的 Python 策略，观察仅截至决策时点的历史，并在下一交易日开盘撮合。
\item 模拟多头、美元账户、整数股、手续费和滑点，记录订单、成交、持仓与权益。
\item 输出净值、收益、风险、成本、换手与基准比较，以及可离线打开的 HTML 报告。
\item 冻结配置、策略、因子、实际输入数据及引擎代码，支持事后复算。
\item 在桌面 App 中选择、编辑、保存策略与因子，设置参数、校验、运行或取消回测，浏览历史结果。
\end{itemize}

\subsection{明确的能力边界}
系统没有券商连接，不会下实盘订单；不支持做空、杠杆、分钟线、盘口撮合、限价单或退市清算。股票池是固定列表，不是历史时点成分数据库。核心运行与测试可离线完成，但安装依赖、可选 Yahoo 数据下载需要网络或事先准备的缓存。

桌面 App 当前依赖工作目录中的 \path{.venv/bin/python}，不内置完整 Python 运行时。把 App 单独复制到一台新电脑不等于完成安装。代码校验与只读数据接口用于约束可信策略，不是运行不可信 Python 的安全沙箱。

\subsection{建议的阅读顺序}
第一次使用可先阅读第 3、4 节；写策略看第 5、6 节；解释回测结果看第 7、8 节；维护代码、迁移环境与重现旧结果看第 9 至 12 节。所有终端示例默认在项目根目录运行，并已激活其虚拟环境。

\section{架构与目录导航}
\subsection{一条回测如何流经系统}
\begin{center}
\begin{tikzpicture}[node distance=6mm,box/.style={draw=teal!65,fill=pale,rounded corners,minimum width=10.8cm,minimum height=8mm,align=center},>=Stealth]
\node[box] (a) {桌面 App / 命令行：选择配置、策略与因子};
\node[box,below=of a] (b) {配置解析 + 数据加载 + 合同校验};
\node[box,below=of b] (c) {交易日时钟：开盘执行旧指令，收盘产生新决策};
\node[box,below=of c] (d) {账户账本 + 权益曲线 + 指标与基准};
\node[box,below=of d] (e) {运行归档 + 快照 + HTML 报告};
\draw[->] (a)--(b);\draw[->] (b)--(c);\draw[->] (c)--(d);\draw[->] (d)--(e);
\end{tikzpicture}
\end{center}

策略负责回答“希望持有什么”，引擎负责回答“何时、以什么价格、买卖多少、收取多少费用”。因子仅从已有历史中计算数值。将这些职责分开，才能区分投资规则变化与成交模型变化。

\subsection{顶层文件夹}
下表包含当前目录中的业务目录及常见生成目录。带点的目录可能被 Finder 隐藏。
\begin{longtable}{p{0.29\linewidth}p{0.65\linewidth}}
\toprule 目录 & 用途与维护方式\\\midrule\endhead
\path{src/quantlab/} & 可安装的 Python 核心。引擎、数据、指标、报告与入口代码均在此。\\
\path{strategies/} & 用户策略与惰性模板；直接决定交易逻辑，修改应保留规则说明与版本。\\
\path{factors/} & 用户批准的纯计算因子；需要配置显式绑定，不是自动扫描执行的信号库。\\
\path{configs/} & YAML 实验定义；\path{examples/} 放基础示例。复制新配置可保留旧实验条件。\\
\path{data/raw/} & 本地原始/冻结输入，包括合成样例与真实历史 Parquet。勿混淆数据来源。\\
\path{data/processed/}, \path{data/cache/}, \path{data/metadata/} & 数据整理、缓存及来源信息的组织位置；是否读写以配置和具体脚本为准，并非所有目录都参与每次运行。\\
\path{runs/} & 每次回测的独立归档。用于结果查看与复现，不应作为临时缓存随意删除。\\
\path{desktop/} & React/TypeScript 界面及 Tauri 原生壳；内部 \path{src/} 是前端，\path{src-tauri/} 是 Rust。\\
\path{scripts/} & 生成样例、校验数据、干净安装验证、独立算术审计等辅助脚本。\\
\path{tests/} & 离线自动化测试，\path{fixtures/} 保存可手算的测试数据。\\
\path{researchProject/} & 独立研究专题；各自保存代码、数据、结果与报告，不能默认与核心引擎行为相同。\\
\path{validation/} & 以前验收与研究运行的日志。它是历史证据，不代表今天已重新测试。\\
\path{dist/}, \path{build/} & App、wheel 等交付物及构建中间文件。源码改动后需要重新构建。\\
\path{.quantlab-desktop/} & 桌面任务快照与日志；运行中不要清理。取消任务可能留下中间目录。\\
\path{.venv/}, \path{.wheelhouse/} & Python 环境与本地依赖安装包。虚拟环境通常不适合直接跨机器复制。\\
\path{.git/} & 版本历史与远程配置；不是业务数据目录。\\
\path{.pytest_cache/}, \path{.ruff_cache/} & 测试和工具缓存；可重建，不等于回测档案。\\
\bottomrule
\end{longtable}

\subsection{核心模块与入口文件}
\path{data/} 处理交易日历、输入校验与截断历史视图；\path{strategy/} 定义合同、上下文、加载和检查；\path{engine/} 管理时钟、撮合、订单与账本；\path{costs/} 计算费用和滑点；\path{metrics/} 定义统计口径；\path{report/} 生成图表与 HTML；\path{utils/} 提供序列化、路径与日志工具。这些子目录均位于 \path{src/quantlab/}。

\path{config.py} 解析严格配置，\path{cli.py} 提供命令行，\path{archive.py} 生成实验档案，\path{desktop_bridge.py} 为 UI 提供文件、配置、任务与结果操作。\path{types.py} 定义公共类型/异常，\path{__init__.py} 等是包入口。

根目录的 \path{pyproject.toml} 定义包和依赖范围；\path{requirements.lock} 固定已验收依赖；\path{README.md} 提供快速说明；\path{AGENTS.md} 记录项目与策略开发约束；V4 工作文档是原始规格。阶段验收文件和 \path{DESKTOP_PROGRESS.md} 是历史记录；出现与当前代码不一致之处，应先核对代码、测试与具体运行档案。

\chapterpage
\section{安装与首次运行}
\subsection{准备 Python 环境}
项目声明 Python 3.12 或更新版本。已有环境可先检查，不必重复创建：
\begin{lstlisting}[language=bash]
cd /Users/yusa/ChatGPTWorkSpace/QuantLab
source .venv/bin/activate
python --version
python -m pip check
quantlab --help
\end{lstlisting}
新环境的标准流程如下；应在确定使用哪个 Python 解释器后执行：
\begin{lstlisting}[language=bash]
python3.12 -m venv .venv
source .venv/bin/activate
python -m pip install -r requirements.lock
python -m pip install -e '.[dev]'
python -m pytest -q
\end{lstlisting}
可编辑安装使 \path{src/} 的修改在该环境中直接生效。依赖固定文件不是全平台安装包；离线安装还需要对应平台的 wheel。当前虚拟环境可能引用外部 Python runtime，迁移或删除该 runtime 后应重建环境，而非只修改激活脚本。

\subsection{第一个离线示例}
\begin{lstlisting}[language=bash]
python scripts/make_sample_data.py
python scripts/validate_data.py configs/default.yaml
quantlab strategies validate \
  strategies/fixed_weight_buy_and_hold.py \
  --config configs/examples/fixed_weight_buy_and_hold.yaml
quantlab backtest run \
  --config configs/examples/fixed_weight_buy_and_hold.yaml
\end{lstlisting}
生成器使用固定种子构造模拟价格；即使股票代码叫 AAPL、TSLA 或 SPY，也不是其真实历史行情。该例用于验证安装与链路。运行后终端输出归档路径与最终资产；用打印出的实际运行编号替换下面的 \code{RUN_ID}：
\begin{lstlisting}[language=bash]
quantlab report open RUN_ID
\end{lstlisting}
在别的目录打开报告时传入 \code{--runs-dir /absolute/path/to/runs}，也可直接在浏览器打开归档内的 \path{report.html}。

\subsection{使用已有真实历史数据}
GitHub 发布不携带真实行情。新克隆应先生成合成 sample，并在 App 选择 sample 基础配置；以下历史实验需要先准备配置指向的数据文件。

现有周频研究配置指向本地 Yahoo 历史 Parquet，不会因为执行回测而重新下载。先校验再运行：
\begin{lstlisting}[language=bash]
python scripts/validate_data.py \
  configs/weekly_equal_weight_aapl_tsla_2020_2025.yaml
quantlab strategies validate \
  strategies/weekly_equal_weight_aapl_tsla.py \
  --config configs/weekly_equal_weight_aapl_tsla_2020_2025.yaml
quantlab backtest run \
  --config configs/weekly_equal_weight_aapl_tsla_2020_2025.yaml
\end{lstlisting}
这个配置从 2020-01-03 的现金基线开始，以便在 2020-01-06 开盘执行首次指令。把开始日期直接改为首次买入日，会改变可用历史、第一笔交易与收益统计基线。

\section{桌面 App 操作指南}
\subsection{打开与连接}
在 Finder 中打开 \path{dist/QuantLab.app}。若未自动连接，进入“工作区设置”，填写项目根目录的绝对路径。它应包含 \path{pyproject.toml}、\path{src/quantlab/}、策略、配置以及可运行的 \path{.venv/bin/python}。原生 App 无需启动浏览器开发服务器。

左侧导航分为研究工作台、回测记录、数据目录和工作区设置。深色/浅色切换只影响显示。数据目录主要列出现有本地文件，不提供完整的数据采购、在线更新或清洗管理平台。

\subsection{从选择文件到运行实验}
\begin{enumerate}
\item 先选择接近目标实验的配置预设，检查策略文件、股票池、基准和日期。
\item 策略下拉菜单读取 \path{strategies/} 下的 Python 文件。选择后阅读实际代码；切换策略后重新检查参数兼容性。
\item 因子下拉菜单用于打开编辑。策略需要导入它时还必须绑定依赖，并在代码中调用对应函数；仅打开因子不会自动生成交易。
\item 在代码框编辑，点击保存或按 Command+S。多个文件可作为标签保留。新建策略复制项目模板，新建因子只有待实现的空模板。
\item 检查日期、本金、股票池、基准、本地行情路径；展开高级设置查看 JSON 参数、仓位模式和交易成本。
\item 点击校验策略；校验通过后运行回测。运行会保存修改、验证配置并冻结本次输入源码。查看日志以了解进度，必要时取消。
\item 确认结果区的运行编号，再查看净值、指标、成交明细和完整报告。
\end{enumerate}

\subsection{编辑、结果与保存的区别}
保存代码会更新本地 Python 文件，但修改表单不自动覆盖原 YAML；本次参数保存在运行快照中。若希望作为长期配置复用，应另建/编辑 \path{configs/} 中的 YAML。

看到历史曲线不代表当前编辑已被测试。只有一次新的成功运行才对应新的配置与代码。历史列表最多返回最近 100 个条目，成交表最多显示前 2,000 笔；更早记录和完整成交仍在磁盘归档。单条损坏归档会单独显示失败，不阻断其他记录；失败标记优先于成功汇总，且成功记录必须有 complete 元数据。

文件在其他编辑器中被更改后，App 用内容指纹阻止静默覆盖。先复制需保留的草稿，再点击重新读取磁盘文件并确认放弃当前编辑。取消任务不会把未完成输出标记为有效结果；检查运行状态，不把残留目录当作成功回测。

\subsection{开发预览与重新打包}
日常使用已构建 App 不需要安装 Node/Rust；开发或重新打包需要 Node.js（Vite 6 支持的版本）、npm、Rust/Cargo 和 macOS 开发工具。本版本将 Vite 更新到兼容的 6.4.5 修复线、source-map-js 更新到 1.2.2；发布时 npm 审计返回零项已知漏洞。
\begin{lstlisting}[language=bash]
cd desktop
npm ci
npm run desktop
# Build the production macOS app:
npm run bundle
\end{lstlisting}
开发浏览器预览可使用：
\begin{lstlisting}[language=bash]
QUANTLAB_WORKSPACE=/Users/yusa/ChatGPTWorkSpace/QuantLab npm run dev
\end{lstlisting}
该入口仅监听 \code{127.0.0.1:1420}；原生版使用 Tauri IPC 调用 Python。新 App 在 \path{desktop/src-tauri/target/release/bundle/macos/QuantLab.app}。确认测试通过、退出旧 App 后再替换交付目录中的版本。修改源码不会自动更新已经打包的 App。

\chapterpage
\section{如何把策略想法变成可执行实验}
\subsection{先明确规则，再写代码}
口头提出策略时，应提供以下内容：交易股票与观察基准；准确数学规则与阈值；首次建仓方式；信号计算区间；决策和执行时点；节假日处理；仓位与现金分配；手续费与滑点；开始/结束日期；缺失数据处理；允许的参数及不希望自动优化的限制。

例如“每周一等权”仍需明确是开盘成交，还是周一收盘决定次日成交。当前用户批准的 AAPL/TSLA 策略选择前一交易周最后一个交易日收盘生成指令，下一周首个交易日开盘成交。

交付通常包含 \path{strategies/name.py}、\path{configs/name.yaml}，若有新的数值特征再增加 \path{factors/name.py}。需要改变引擎时必须单独识别并测试，不能为改善收益暗中改变规则。表达清楚策略后，AI 也应遵循同一过程：写规则说明、实现、校验、回测、核查和报告。

\subsection{自己写策略的接口}
从 \path{strategies/_template.py} 复制起步，或在 App 中新建。下面是没有交易逻辑的最小骨架，不是额外推荐的研究策略：
\begin{lstlisting}[language=Python]
from quantlab.strategy.base import BaseStrategy

STRATEGY_META = {
    "name": "my_strategy", "version": "0.1.0",
    "description": "Document the exact approved rule",
}

class Strategy(BaseStrategy):
    def warmup_bars(self):
        return 0

    def generate_target_weights(self, dt, data, ctx):
        return None

    def decision_metadata(self):
        return {}

def build_strategy(params):
    return Strategy(params)
\end{lstlisting}
\code{build_strategy(params)} 创建实例；基类复制参数，策略可保存其自身的历史状态。\code{decision_metadata()} 用于返回可审计的决策解释，例如因子数值和触发原因，不应另行下单。

返回值语义必须区分：\code{None} 表示没有新指令；\texttt{\{\}} 表示目标全现金；字典表示整个目标组合，没有列出的现有持仓目标为零。权重必须有限、非负、属于可交易股票池，且不超过单股/总敞口限制。若权重之和小于 1，剩余部分对应现金目标。价格漂移后不会自动强制恢复权重。

\subsection{可以读取什么}
\begin{longtable}{p{0.43\linewidth}p{0.51\linewidth}}
\toprule 接口 & 语义\\\midrule\endhead
\code{data.history(symbol, bars=None, end=None)} & 只返回决策时点之前可见的历史副本，可限制条数和结束时间。\\
\code{data.current(symbol)} & 当前可见 bar 或 \code{None}，应处理缺失。\\
\code{data.as_of} & 数据视图的截止交易日。\\
\code{ctx.cash}, \code{ctx.equity} & 当前只读现金和权益。\\
\code{ctx.positions}, \code{ctx.universe} & 只读股数映射与允许交易的股票池。\\
\code{ctx.next_session_date} & 已知交易日历的下一 session 日期，不是未来价格。\\
\code{ctx.seed}, \code{ctx.execution_sizing} & 随机种子及实际仓位转换模式。\\
\bottomrule
\end{longtable}
基准历史可用于观察，但不能因此把基准当作可交易股票。\code{warmup_bars()} 只累计回测区间内的可观测 bar；没有自动加载区间外训练历史，需在配置中预留所需数据。

策略不应下载行情、读取未来数据、调用引擎私有对象或直接修改账本。校验器会检查常见违规导入和文件 I/O，但不能证明任意 Python 程序安全，也不能证明投资逻辑有效。

\subsection{接入配置与运行}
复制现有 YAML，修改 \code{strategy.path}、\code{params}、股票池和数据路径；相对路径以 YAML 所在目录为基准。CLI 的显式 \code{--strategy} 路径则相对于当前终端目录，并覆盖 YAML 指定策略。
\begin{lstlisting}[language=bash]
quantlab strategies validate strategies/my_strategy.py \
  --config configs/my_strategy.yaml
quantlab backtest run --config configs/my_strategy.yaml
\end{lstlisting}
校验通过只意味着基本合同满足。仍需检查极端输入、关键交易日、仓位与现金，以及输出是否真正符合书面规则。

\section{因子库与已有策略}
\subsection{因子与策略不是同一件事}
因子是计算出来的特征；策略是使用特征、时间和账户状态决定交易的规则。例如每日 AAPL 相对 SPY 的收益差为
\[
 f_t=\left(\frac{P^{A}_{t}}{P^{A}_{t-1}}-1\right)
     -\left(\frac{P^{M}_{t}}{P^{M}_{t-1}}-1\right).
\]
这里使用系统提供的复权价格。该差值不是 beta 调整后的残差，也不是对个股“高估/低估”的直接证明。以 \(\pm0.01\) 为阈值是相差正负 1 个百分点，不是相对收益差再乘百分比。

因子不决定买入多少，也不决定何时交易。同一个因子可以被不同策略使用；反过来，“每周恢复 50/50”只是固定权重再平衡，不需要收益预测因子，所以不会自然产生因子库文件。

\subsection{当前因子的实际接口}
\path{factors/return_spread.py} 提供 \code{daily_return_spread} 和 \code{weekly_return_spread}。前者使用同步的两日收盘，后者使用两个相邻完整交易周的最后收盘。周收益不能用最后一天的单日收益替代，也不简单等于一周日收益的算术和。

周函数应在策略已确认本周最后一个交易日的收盘时调用，否则可能把未完成周当作完整周。历史不足返回 \code{None}；端点不同步、非正价格等会报错。函数返回起止日期、个股收益、基准收益和差值，用于审计。

绑定示例（位于 \path{configs/} 中）：
\begin{lstlisting}
strategy:
  path: ../strategies/aapl_weekly_relative_reversal.py
  params:
    start_date: '2020-01-06'
    end_date: '2025-12-31'
    threshold: 0.01
  rebalance_schedule: daily
  factor_paths:
    factors.return_spread: ../factors/return_spread.py
execution:
  sizing: decision_relative
\end{lstlisting}
这是配置片段，其他部分应保留对应完整配置。策略中通过 \code{from factors.return_spread import weekly_return_spread} 调用。加载器按模块名显式绑定、检查并冻结源码，不自动执行因子目录内所有文件。

\subsection{已有策略的区别}
\begin{longtable}{p{0.43\linewidth}p{0.51\linewidth}}
\toprule 文件 / 实验 & 实际规则\\\midrule\endhead
\path{fixed_weight_buy_and_hold.py} & 初次发出给定权重，随后不主动再平衡。基础示例各半 AAPL/TSLA。半仓 AAPL 对照实验复用此实现。\\
\path{weekly_equal_weight_aapl_tsla.py} & 从 2020-01-06 开始，下周首个交易日开盘恢复各 50\% 目标；无需因子。\\
\path{aapl_weekly_relative_reversal.py} & 初始一半资金买 AAPL；之后周收益差低于 -1pp 买半数现金，高于 +1pp 卖半数股数，否则保持。\\
\path{_template.py} & 惰性起点，返回 \code{None}，不代表有效投资规则。\\
\bottomrule
\end{longtable}
相对反转策略首次建仓优先，不在同一时点叠加因子交易；现金不生息，买入预算含手续费，卖出向下取整。阈值相等（代码含 \(10^{-12}\) 的边界容差）不交易。它要求 \code{decision_relative} 模式，不能与普通目标权重转换随意互换。

\chapterpage
\section{数据、配置与执行模型}
\subsection{数据合同与复权}
本地输入为 long format，每行一只股票一个交易日：
\begin{lstlisting}
date,symbol,open,high,low,close,volume,adj_close
\end{lstlisting}
\code{date} 是无时区的交易日标签，系统决策与成交时间戳使用 UTC。数据需要覆盖股票池和基准，并符合 XNYS 交易日历。重复日期/股票、非法价格、负成交量、非交易日和缺少所需股票会明确失败。

在 \code{adjusted} 模式下，使用 \(a_t=\text{adj\_close}_t/\text{close}_t\) 同比缩放 OHLC；不再单独入账拆股和分红，成交量保持源值。整数股是复权研究单位，不应解释为历史原始股数。\code{raw} 模式不处理公司行动，显式非零分红或非平凡拆股标记会被拒绝；漏报行动仍会产生偏差。

更新数据时应保留旧文件，写入新文件并记录来源、下载日期、区间、复权政策和校验结果。复权历史可能被源端修订；固定旧输入才能解释两次结果差异。Yahoo adapter 是可选便利入口，不是唯一权威数据源。需要时先安装 \code{'.[online]'}，显式选择 \code{data.provider: yahoo}；桌面主流程仍以本地文件为主。

\subsection{主要配置项}
\begin{longtable}{p{0.31\linewidth}p{0.63\linewidth}}
\toprule 配置段 & 控制内容\\\midrule\endhead
\code{backtest} & 起止日期、本金（正数）、USD 基准货币。\\
\code{data} & provider、path、1d 频率、复权方式、基准、来源描述和异常收益告警阈值。\\
\code{universe} & 非空、无重复的可交易代码列表，以及固定列表/当前成分来源声明。\\
\code{strategy} & 源码、参数、决策资格频率和因子模块到源码路径的映射。\\
\code{execution} & 收盘决策、下一开盘执行、整数股及 sizing 模式。\\
\code{commission}, \code{slippage} & 手续费模型、最低费和基点滑点。1 bps 等于 0.01\%。\\
\code{portfolio} & 禁止做空和负现金，单股权重与总敞口限制。\\
\code{analytics} & 年化交易日数（默认 252）、无风险率、滚动窗口。\\
\code{report} & 报告开关、是否自动打开、运行输出目录。\\
\code{reproducibility} & 随机种子。仅固定种子不能替代数据和代码快照。\\
\bottomrule
\end{longtable}
Pydantic 严格校验未知字段和非有限数值；YAML 重复键也会报错。配置路径相对于 YAML，而不是 App 当前界面位置。

\subsection{日内事件顺序与周频的含义}
\begin{enumerate}
\item 当日开盘处理上一交易日收盘形成的待执行目标。
\item 按当日开盘标记已有仓位，冻结目标转换所需权益。
\item 结合滑点和整数股计算订单，先卖后买，买入按代码字母顺序处理。
\item 现金不足时确定性缩量并取消余量，不允许借款；零成交不收最低手续费。
\item 当日收盘估值、生成截断历史视图、调用策略、记录日终状态。
\end{enumerate}
\code{rebalance_schedule: weekly} 指“每周首个交易日收盘具有决策资格”，不等于“周一开盘成交”。已有周一策略使用 \code{daily} 每天检查下一 session 是否进入新周，实际只在周末边界发出目标。已知交易日历可以提前使用，未来价格不能。

\subsection{两种仓位转换方式}
普通 \code{target_weights} 按开盘权益和目标权重转换成整数数量。即使权重写为各半，费用、滑点、整数舍入与现金不足都可能使实际权重偏离 1:1。

\code{decision_relative} 在收盘把权重变化转换并冻结为现金/股数比例，下一开盘按该比例执行。这样周末跳空后，“用一半剩余现金买入”仍保持原意。此模式目前仅支持单只可交易股票，用于已有 AAPL 相对反转；不是多资产通用仓位算法。

\subsection{成本与缺失数据}
手续费支持零费、每股费加最低费、成交额基点费；滑点支持零或基点模式，买入加价、卖出减价。滑点已包含在成交价中，不能再重复扣现金。买入手续费进入平均成本，卖出净收入计入已实现损益。

缺少执行 bar、开盘价无效或成交量为零时取消订单；不得借未来价格补成交。持仓缺少当日 bar 时用最后已知价格估值并告警；缺少基准 bar 则失败。最终一日的待执行指令取消，但期末仓位不会强制卖出。

日线最终成交量只用于零量可交易标记；真实开盘时尚不知道当日总量，所以这是模型假设，不是盘中流动性模拟。账本持续检查
\[
E_t=C_t+MV_t=E_0+\mathrm{RealizedPnL}_t+\mathrm{UnrealizedPnL}_t.
\]

\section{如何阅读结果而不误读指标}
先核对数据来源、开始/结束日期、首次买入日、基准、成本和警告，再看收益。较低回撤可能来自较低股票敞口，而非更好的选时；相对反转应与匹配的初始半仓买入持有对照比较，也可同时观察 SPY。

令 \(E_t\) 为每日收盘权益，\(r_t=E_t/E_{t-1}-1\)。首个 session 收盘为现金基线，日收益统计排除第一行占位零收益。核心口径为：
\begin{align*}
R_{\rm total}&=E_T/E_0-1,\\
\mathrm{CAGR}&=(E_T/E_0)^{1/Y}-1,\quad Y=\frac{d_T-d_0}{365.25},\\
\sigma_{\rm ann}&=s(r_t)\sqrt{252},\\
\mathrm{Sharpe}&=\frac{\overline{r_t-r_f}}{s(r_t-r_f)}\sqrt{252},\\
\mathrm{MDD}&=\min_t\left(\frac{E_t}{\max_{u\le t}E_u}-1\right).
\end{align*}
252 可由配置改变；\(s\) 使用样本标准差，年化无风险率按几何方式转成日率。Sortino 分母为所有观测的负超额收益平方均值的平方根；Calmar 为 CAGR 除以最大回撤绝对值。回撤持续时间按日历日计算，未修复回撤截至最后观察日。

换手为双边实际成交额除以前一收盘权益，再按收益区间数年化，不除以二。平均现金与敞口包含初始现金日。Beta、alpha、tracking error 和 information ratio 使用与策略对齐的基准日收益；基准是首日收盘开始、无费用且可分数股的买入持有曲线，不等同于策略的真实成交账户。

无法定义的比率为 \code{null}/N/A，不应替换为零；成交笔数不是已完成往返交易次数。系统不为买入持有强行给出往返胜率。HTML 报告内嵌图片，可离线阅读；数值审计应以 CSV/Parquet/JSON 原始输出为依据。

\chapterpage
\section{归档、复现与独立审计}
\subsection{每次运行保存什么}
运行目录名包含 UTC 时间、策略名和摘要，重复运行不会覆盖旧目录。下列文件用于回答“当时到底运行了什么”：
\begin{longtable}{p{0.43\linewidth}p{0.51\linewidth}}
\toprule 文件 & 作用\\\midrule\endhead
\path{config.snapshot.yaml} & 指向冻结输入的解析配置。\\
\path{config.original.json} & 初始解析配置，用于追溯原路径。\\
\path{strategy.snapshot.py}, \path{factors.snapshot/} & 策略与显式依赖因子源码。\\
\path{data.snapshot.parquet} & 当次使用的调整前输入 panel，避免重新联网。\\
\path{code.snapshot/}, \path{requirements.snapshot.txt} & 当时的核心代码、安装描述和依赖版本。\\
\path{metadata.json}, \path{manifest.json} & 环境、版本、Git 状态、指纹与文件 SHA-256。清单不含自身及持续写入的日志。\\
\path{equity.csv}, \path{returns.csv}, \path{benchmark.csv} & 权益、日收益和基准。\\
\path{decisions.csv}, \path{orders.csv}, \path{fills.csv} & 决策原因、请求/取消/成交及成本。\\
\path{positions.parquet} & 逐日持仓。\\
\path{summary.json}, \path{metrics.json}, \path{warnings.json} & 汇总、完整指标与告警。\\
\path{report.html}, \path{logs.txt}, \path{reproduce.sh} & 人类可读报告、运行日志和复现入口。\\
\path{failure.json} & 失败标记；不能把该目录当作成功实验。\\
\bottomrule
\end{longtable}
报告关闭时可以没有 HTML。源码有未提交修改时，Git commit 本身不足以确定运行内容；冻结源码和哈希更直接。不要只备份截图或 PDF 而丢失输入和账本。

\subsection{重跑与迁移}
同一环境中运行：
\begin{lstlisting}[language=bash]
sh runs/RUN_ID/reproduce.sh
\end{lstlisting}
这会生成新的运行目录，不改写旧结果。跨环境应先创建新的虚拟环境，安装该次 \path{requirements.snapshot.txt} 与 \path{code.snapshot/}，再执行相应复现命令。归档移动后，配置中的绝对路径可能失效；保留原件，另存修改路径后的配置，并记录修改。不要用当前源码冒充历史源码复算。

同输入、同规则、同环境应得到一致经济结果；跨平台浮点差异不保证文件逐字节相同。对比时同时检查交易日期、股数、费用、每日权益和指标，而不是只看最终收益。

\subsection{现有策略的独立检查}
\begin{lstlisting}[language=bash]
python scripts/audit_weekly_relative_reversal.py runs/RUN_ID
\end{lstlisting}
该脚本独立重算现有周频反转规则的信号、现金/股数与每日权益，不调用策略和引擎实现。它适用于对应规则和成本条件，不是任意策略的通用审计器；更改规则后需相应更新独立计算依据。

\section{现有研究专题如何使用}
\path{researchProject/RESEARCH_INDEX.md} 是研究入口。当前四个专题分别研究：
\begin{enumerate}
\item \path{01_aapl_tsla_rebalancing/}：AAPL/TSLA 目标权重和日历再平衡频率。
\item \path{02_threshold_partial_rebalancing/}：阈值触发与部分纠偏对交易频率、跟踪和表现的影响。
\item \path{03_time_and_asset_robustness/}：入场时间、持有期限和资产组合的稳健性。
\item \path{04_rebalancing_mechanisms/}：分散化增长与相对买入持有净超额收益的区别。
\end{enumerate}
每个专题按自身 README/审计说明复算，阅读其报告、数据字典、协议、结果与代码，而不是直接在 App 选择报告文件。后续专题有独立计算引擎；不能因名称相近就假设与主引擎逐项相同。

尤其第 04 个专题区分无摩擦收盘到收盘恒等式、合成实验和原有开盘成交约定。数学分解项不能直接替代含成本的开盘回测收益。固定资产选择、重叠窗口和参数比较也不能自动构成样本外证据。要把研究发现转成正式策略，应重新明确规则、写入策略/因子合同、增加测试并验证主引擎结果。

\chapterpage
\section{维护、测试与交付}
\subsection{推荐的修改流程}
\begin{enumerate}
\item 先写变更目的与预期行为，确认修改属于策略、数据、引擎还是界面。
\item 保存原配置与参考运行，使用版本控制记录源码；不要覆盖旧行情文件。
\item 实现最小必要改动。引擎改动必须增加能发现真实错误的测试，保留未来价格扰动等防前视检查。
\item 运行相关测试和完整套件；必要时做一条本地离线端到端回测。
\item 对有交易语义影响的改动，比较订单、成交、权益和成本，说明结果变化原因。
\item 更新手册/README 和配置说明；界面改动后重新构建并验证实际 App。
\end{enumerate}

\subsection{日常检查命令}
\begin{lstlisting}[language=bash]
python -m pytest -q
python -m pip check
# Desktop checks, from desktop/:
npm ci
npm run build
cargo check --manifest-path src-tauri/Cargo.toml
npm run bundle
\end{lstlisting}
本次文档编写于 2026-10-22 重新运行了 Python 完整测试：\textbf{136 项通过}。这不表示所有后续改动都会通过，也不替代 App 发布前的原生窗口检查。测试覆盖账本、费用、缺失数据、策略校验、防前视、配置、归档复现与桌面桥接。

\path{scripts/check_clean_install.py} 可在临时环境验证 wheel 安装、测试和离线样例；前提是 \path{dist/} 中已有项目 wheel，\path{.wheelhouse/} 中有全部所需平台依赖。没有这些准备时，不应把脚本失败误判为策略失败。

\subsection{依赖、备份与发布边界}
依赖升级应在新环境中做，核对 \path{pyproject.toml} 的兼容范围，更新相应锁文件并跑回归。前端使用 \path{package-lock.json}，Rust 使用 \path{Cargo.lock}；不要只升级一端而忽略原生接口兼容性。

建议备份源码、配置、原始数据、完整重要运行档案与研究协议。环境和构建缓存可重建，但若要求完全离线恢复，还需备份兼容的 wheel 及安装工具。清理前先确认没有运行中任务；不要把 \path{runs/} 与 \path{.quantlab-desktop/} 当成同一种目录。

2026-10-22 用户批准将现有简单策略和因子一并发布。2026.10.22 版本包含本体、桌面源码、配置、测试及本手册；不包含独立 \path{researchProject/}、真实行情、运行档案、环境与缓存。现有 GitHub 仓库在发布核对时为公开状态；本次保留其可见性。维护时核对发布清单与 Git 历史；\path{.gitignore} 不会移除已提交内容。发布边界见根目录 \path{PROJECT_SCOPE.md}。

\subsection{手册本身的维护}
本文件是自包含 LaTeX 源码，无外部图片依赖。中文使用 XeLaTeX 和系统宋体/黑体（或 Fandol 字体回退）；英文使用 TeX Gyre 字体。修改后编译两次以更新目录与交叉引用，查看 PDF 是否有溢出、缺字或断裂代码。用临时输出目录保存辅助文件，最终 cn/en 目录只保留源文件和 PDF。

\section{故障排查与操作检查表}
\begin{longtable}{p{0.27\linewidth}p{0.67\linewidth}}
\toprule 现象 & 先检查什么\\\midrule\endhead
命令不存在 / 包无法导入 & 当前目录与虚拟环境；\code{python --version}、可编辑安装、\code{quantlab --help}。\\
App 连接失败 & 工作区是否是项目根目录；\path{.venv/bin/python} 能否运行；目录是否移动、依赖是否完整。\\
配置不合法 & 未知/重复 YAML 字段、日期顺序、参数类型；相对路径基于 YAML。\\
策略校验失败 & 元信息、基类、构造函数、返回权重、违规导入、显式因子绑定和 sizing 模式。\\
缺少因子模块 & \code{factor_paths} 的模块名是否与 import 一致；是否确实绑定该源码。\\
没有成交 & 是否返回 \code{None}、预热不足、日期未到、资金不足，或订单因无开盘/零成交量被取消。看 decisions/orders。\\
周频晚一天 & 是否误把 weekly 决策资格当作周一开盘执行；检查下一交易日日期。\\
权重不是精确各半 & 整数股、成本、开盘跳空和现金约束；区分目标和实际权重。\\
曲线没有变化 & 当前显示是否为历史运行；修改代码/表单后是否产生新的成功编号。\\
保存发生冲突 & 外部文件已更改；保留草稿，再重载并人工合并，勿强行覆盖。\\
运行失败 / 报告为空 & 查 \path{failure.json}、日志、报告开关；不要仅凭目录存在判断成功。\\
复现结果不同 & 数据/代码/参数/依赖/种子是否一致；是否误用新引擎；归档路径是否失效。\\
更新源码但 App 不变 & 是否重新运行 bundle，并打开新生成的 App 而非旧副本。\\
\bottomrule
\end{longtable}

每次研究提交前核对：规则是否事先确定；数据是否真实且来源可追溯；时点是否无前视；成本是否明确；异常是否披露；是否有合适对照；是否保存完整运行档案。收益不佳应如实报告，不能自动修改阈值直到回测好看。

\appendix
\section{代码依据与版本说明}
本手册于 2026-10-22 核对以下本地材料：\path{README.md}、\path{AGENTS.md}、\path{pyproject.toml}、\path{configs/}；核心 \path{config.py}、\path{cli.py}、策略基类/上下文、因子函数、绩效实现、桌面桥接；\path{desktop/README.md}；\path{researchProject/RESEARCH_INDEX.md} 与辅助脚本。本次本体版本标识为 \code{2026.10.22}，对应 Git 标签 \code{v2026.10.22}；具体提交由仓库标签解析。审查范围和证据见 \path{AUDIT_2026.10.22.md}。

本文以实际实现为准。旧 README 的 V2 路线图和“唯一官方示例”描述对应较早阶段；当前已经存在用户批准的追加策略、桌面界面及研究专题。手册不会据此推断尚未实现的实时交易、通用自动寻优、完整因子评估平台或跨机器独立安装能力。
\end{document}
